\documentclass[11pt,a4paper]{article}

\usepackage[margin=1in]{geometry}
\usepackage{amsmath,amssymb,amsthm}
\usepackage{mathtools}
\usepackage{booktabs}
\usepackage{microtype}
\usepackage{hyperref}
\usepackage{listings}
\usepackage{xcolor}
\usepackage{tikz}
\usetikzlibrary{cd,arrows,shapes}

\hypersetup{
    colorlinks=true,
    linkcolor=blue!70!black,
    citecolor=green!50!black,
    urlcolor=blue!80!black
}

\lstdefinelanguage{Lean}{
  keywords={structure, extends, def, noncomputable, theorem, lemma, inductive, where, let, in, match, with, fun, if, then, else, Prop, Type},
  sensitive=true
}

\lstset{
  language=Lean,
  basicstyle=\ttfamily\small,
  columns=fullflexible,
  breaklines=true,
  commentstyle=\color{gray},
  keywordstyle=\color{blue!80!black}\bfseries,
  stringstyle=\color{teal},
  frame=single,
  frameround=tttt,
  backgroundcolor=\color{gray!5},
  extendedchars=true,
  literate=
    {->}{{$\to$}}1
    {<-}{{$\leftarrow$}}1
    {∀}{{$\forall$}}1
    {∃}{{$\exists$}}1
    {∧}{{$\land$}}1
    {∨}{{$\lor$}}1
    {ω}{{$\omega$}}1
    {δ}{{$\delta$}}1
    {μ}{{$\mu$}}1
    {σ}{{$\sigma$}}1
    {τ}{{$\tau$}}1
    {φ}{{$\phi$}}1
    {Φ}{{$\Phi$}}1
    {ψ}{{$\psi$}}1
    {⪯}{{$\preceq$}}1
    {≃ᵣ}{{$\simeq_r$}}1
    {↦}{{$\mapsto$}}1
    {Γ}{{$\Gamma$}}1
    {Δ}{{$\Delta$}}1
    {α}{{$\alpha$}}1
    {β}{{$\beta$}}1
    {↔}{{$\leftrightarrow$}}1
}

\newtheorem{theorem}{Theorem}[section]
\newtheorem{lemma}[theorem]{Lemma}
\newtheorem{proposition}[theorem]{Proposition}
\newtheorem{corollary}[theorem]{Corollary}
\theoremstyle{definition}
\newtheorem{definition}[theorem]{Definition}
\newtheorem{example}[theorem]{Example}
\theoremstyle{remark}
\newtheorem{remark}[theorem]{Remark}

\newcommand{\HAomega}{\ensuremath{\mathrm{HA}^\omega}}
\newcommand{\Rep}{\ensuremath{\mathbf{Rep}}}
\newcommand{\Q}{\ensuremath{\mathbb{Q}}}
\newcommand{\N}{\ensuremath{\mathbb{N}}}
\newcommand{\R}{\ensuremath{\mathbb{R}}}
\newcommand{\EFTC}{\ensuremath{\mathrm{EFTC}}}
\newcommand{\Th}{\ensuremath{\operatorname{Th}}}
\newcommand{\Req}{\ensuremath{\operatorname{Req}}}
\newcommand{\mr}{\ensuremath{\;\mathbf{mr}\;}}
\newcommand{\preceqR}{\ensuremath{\preceq}}
\newcommand{\simeqR}{\ensuremath{\simeq_r}}

\title{\textbf{How Much Data Does a Theorem Need?\\Representation Adequacy for Constructive Analysis, Mechanized in Lean 4}}

\author{Ara Aslyan}
\date{\today}

\begin{document}

\maketitle

\begin{abstract}
We present a mechanized calibration framework for constructive real analysis formalized in the Lean~4 interactive theorem prover. Rather than asking whether real analysis is constructive, we investigate a quantitative, fine-grained reverse-mathematics question: \emph{how much computational data must a representation of a real function carry before a given classical theorem becomes provable about it?}

By formalizing a ladder of function representations---from continuous functions with a modulus of uniform continuity ($A_0$), to continuously differentiable functions with derivative modulus and derivative code ($A_1$), and their approximating counterparts ($E_0, E_1$)---we prove that fundamental theorems of calculus sort themselves along this hierarchy:
\begin{enumerate}
    \item \textbf{The First Fundamental Theorem of Calculus ($\text{EFTC1}$)} holds at $A_0$ with only a modulus of uniform continuity ($\delta := \omega$), requiring no derivative data.
    \item \textbf{The Second Fundamental Theorem of Calculus ($\text{EFTC2}$)} requires the enriched representation $A_1$.
    \item \textbf{The bridge between them costs exactly \texttt{Classical.choice}:} the coercion \texttt{A0.toA1} from a uniformly differentiable function in $A_0$ to $A_1$ is noncomputable with an axiom footprint of \emph{precisely} \texttt{[Classical.choice]}, providing a mechanical witness to Myhill's computability obstruction.
\end{enumerate}
We further structure this representation space into a category $\Rep$ equipped with a retract preorder ($\preceq$) and prove constructive cross-representation Galois equivalences ($A_1 \simeqR A_0^{\text{diff}}$ and $E_1 \simeqR E_0^{\text{diff}}$). All theorems are fully verified in Lean~4 across 7,895 compilation jobs without non-constructive artifacts in core extraction, establishing kernel \texttt{\#print axioms} auditing as a quantitative measurement instrument.
\end{abstract}

\section{Introduction}

Constructive analysis, from Brouwer and Bishop~\cite{bishop1967,bishop1985} to modern Type-Two Effectivity (TTE)~\cite{weihrauch2000}, has established that classical analysis theorems diverge in their algorithmic content. In standard constructive formalizations, however, theorems are typically either proved under global constructive hypotheses or set aside as non-constructive.

In this work, we develop a mechanized calibration framework in Lean~4~\cite{demoura2021} that formalizes the exact data requirements of real analysis theorems. We investigate the central question:
\begin{quote}
\centering
\textbf{How much data must a representation of a real function carry before a given classical theorem becomes provable about it?}
\end{quote}

The answer is organized as a structured \emph{ladder of representations}, summarized by three foundational results verified in our development:
\begin{enumerate}
    \item \textbf{Integration is Cheap ($\text{EFTC1}$ at $A_0$):}
    \[
    \int_a^b f'(t) \, dt = f(b) - f(a)
    \]
    The First Fundamental Theorem of Calculus (\texttt{eftc1}) holds constructively for functions in $A_0$. The required modulus of integration is supplied directly by the modulus of uniform continuity ($\delta := \omega$). No derivative data or differentiability modulus is needed.
    \item \textbf{Differentiation is Expensive ($\text{EFTC2}$ at $A_1$):}
    \[
    \frac{d}{dx} \left( \int_a^x f(t) \, dt \right) = f(x)
    \]
    The Second Fundamental Theorem of Calculus (\texttt{eftc2\_thm}) requires $A_1$, which carries both explicit derivative codes ($f'$) and a modulus of uniform differentiability ($\delta$).
    \item \textbf{The Separation Gap is Exactly \texttt{Classical.choice}:}
    The coercion \texttt{A0.toA1 (A : A0) (h : HasUnifDeriv A) : A1} is \texttt{noncomputable} with an axiom footprint of \textbf{exactly \texttt{[Classical.choice]}}. While a uniform derivative is mathematically guaranteed to exist by $h$, constructing the representation $A_1$ requires non-constructively choosing the moduli witnesses.
\end{enumerate}

\subsection{The Unifying Principle: The Modulus is the Content}
Across uniform limits, Riemann integration, functional inversion, and differential equations, our formalization validates a single unifying principle:
\begin{quote}
\centering
\textbf{The object is never in doubt. The modulus is the algorithmic content.}
\end{quote}
Classically, the derivative, the Riemann integral, and the inverse function exist by compactness and completeness. Constructively, what distinguishes tractable theorems from non-constructive principles is whether the representation provides the quantitative moduli ($\omega, \delta, \mu$) necessary to compute approximations to arbitrary precision $\varepsilon = 2^{-k}$.

\subsection{Methodological Contribution: \texttt{\#print axioms} as a Quantitative Instrument}
We introduce a proof-theoretic methodology: \textbf{quantitative axiom auditing via the Lean~4 kernel}. In an interactive development whose core term language and natural deduction calculus are choice-free by construction, the axiom footprint reported by Lean's trusted kernel (\texttt{\#print axioms}) serves as a precise measuring instrument. When \texttt{A0.toA1} reports \texttt{[Classical.choice]}, it mechanically certifies that the step from ``a derivative exists'' to ``here is its modulus and evaluator'' requires non-constructive choice, exhibiting Myhill's theorem~\cite{myhill1971} directly inside the proof assistant.

\section{The Representation Ladder}

To capture the analytical distinction between exact point evaluators and numerical approximations, we formalize four foundational representations of real functions on compact intervals $[a, b] \subseteq \Q$:

\begin{figure}[htbp]
\centering
\begin{tikzcd}[row sep=large, column sep=large]
  \begin{array}{c} \mathbf{RepA1} \\ \langle f, \omega, f', \delta, \text{diff}\rangle \end{array} \arrow[d, "\text{forget derivative}"'] \arrow[r, "\text{inclusion}"] &
  \begin{array}{c} \mathbf{RepE1} \\ \langle f, \omega, f', \delta, \text{diff}, \text{Cauchy}\rangle \end{array} \arrow[d, "\text{forget derivative}"] \\
  \begin{array}{c} \mathbf{RepA0} \\ \langle f, \omega, \text{cont}\rangle \end{array} \arrow[r, "\text{inclusion}"'] \arrow[ur, "\text{integration}"] &
  \begin{array}{c} \mathbf{RepE0} \\ \langle f, \omega, \text{cont}, \text{Cauchy}\rangle \end{array}
\end{tikzcd}
\caption{The Smoothness and Approximation Ladder.}
\end{figure}

\subsection{Exact Rational Evaluators: $A_0$ and $A_1$}
\begin{definition}[Representation $A_0$]
A function in $A_0$ consists of an exact rational evaluator $f : \Q \to \Q$ together with a modulus of uniform continuity $\omega : \N \to \N$:
\begin{lstlisting}[language=Lean]
structure A0 where
  a b : Q
  f   : Q -> Q
  ω   : Nat -> Nat
  cont : ∀ k x y, |x - y| < 2^(-ω k) -> |f x - f y| < 2^(-k)
\end{lstlisting}
\end{definition}

\begin{definition}[Representation $A_1$]
A function in $A_1$ extends $A_0$ with an explicit derivative code $f' : \Q \to \Q$ and a modulus of uniform differentiability $\delta : \N \to \N$:
\begin{lstlisting}[language=Lean]
structure A1 extends A0 where
  f'   : Q -> Q
  δ    : Nat -> Nat
  diff : ∀ k x h, 0 < |h| ∧ |h| < 2^(-δ k) -> 
           |(f (x + h) - f x) / h - f' x| < 2^(-k)
\end{lstlisting}
\end{definition}

\subsection{Why the Approximating Layer ($E_0, E_1$) is Necessary}
A fundamental mathematical observation motivates the $E$ layer: \textbf{the Riemann integral of an exact rational-valued function is not rational-valued in general.}

Even if $f \in A_0$ evaluates to exact rationals on $\Q$, its integral $\int_a^x f(t) \, dt$ requires infinite Cauchy sequences of Riemann sums. Thus, integration is an operation that necessarily leaves the $A$ layer and lands in the approximating layer $E_0$:
\begin{lstlisting}[language=Lean]
structure E0 where
  a b  : Q
  f    : Q -> Nat -> Q
  ω    : Nat -> Nat
  cauchy : ∀ k1 k2 x, |f x k1 - f x k2| < 2^(-k1) + 2^(-k2)
  cont   : ∀ k x y, |x - y| < 2^(-ω k) -> |f x k - f y k| < 2^(-k)
\end{lstlisting}

\section{Calibration Results}

Every theorem in Table~\ref{tab:calibration} has been mechanized in the repository and verified via \texttt{\#print axioms}.

\begin{table}[htbp]
\centering
\small
\begin{tabular}{lll}
\toprule
\textbf{Theorem / Construction} & \textbf{Statement in Lean 4 Development} & \textbf{Axiom Footprint} \\
\midrule
EFTC1 at $A_0$ & \texttt{eftc1 (A : A0) : EFTC1Claim A} & \texttt{[propext, choice, Quot]} \\
EFTC2 at $A_1$ & \texttt{eftc2\_thm (A : A1) : EFTC2Claim A} & \texttt{[propext, choice, Quot]} \\
Myhill Choice Bridge & \texttt{A0.toA1 (A : A0) (h : HasUnifDeriv)} & \textbf{\texttt{[Classical.choice]}} \\
EFTC2 from Uniform Deriv & \texttt{eftc2\_of\_unifDeriv (A : A0) h} & \texttt{[propext, choice, Quot]} \\
Integral Lands in $E_0$ & \texttt{A0.intE0 (A : A0) : E0} & \texttt{[propext, choice, Quot]} \\
Integral Lands in $E_1$ & \texttt{A0.intE1 (A : A0) : E1} & \texttt{[propext, choice, Quot]} \\
$A_0$ Composition Closure & \texttt{CompData.comp (F G : A0) : A0} & \texttt{[propext, choice, Quot]} \\
Uniform Limit Closure & \texttt{LimSeq.toE0 (S : LimSeq) : E0} & \texttt{[propext, choice, Quot]} \\
Monotone Inversion & \texttt{inv\_modulus (A : A0) ($\mu$ : Nat -> Nat)} & \texttt{[propext, choice, Quot]} \\
\bottomrule
\end{tabular}
\caption{Mechanized Calibration Theorems and Measured Axiom Footprints.}
\label{tab:calibration}
\end{table}

\subsection{EFTC1 Holds at $A_0$ with $\delta := \omega$}
We prove that Riemann integration requires only uniform continuity:
\[
\left| \sum_{i=0}^{n-1} f'(\xi_i) \Delta x_i - (f(b) - f(a)) \right| < 2^{-k}
\]
In \texttt{eftc1}, the partition modulus $\delta(k)$ is set directly to $\omega(k+1)$. The proof proceeds constructively without requiring higher derivative bounds.

\subsection{EFTC2 Requires $A_1$}
Differentiating the integral requires establishing that:
\[
\lim_{h \to 0} \frac{1}{h} \int_x^{x+h} f(t) \, dt = f(x)
\]
In \texttt{eftc2\_thm}, the differentiability modulus $\delta(k)$ of the integral is explicitly supplied by $f$'s continuity modulus $\omega(k+1)$, and the derivative code is identified with $f$ itself.

\subsection{The Myhill Separation and Classical Choice}
To transition from $A_0$ to $A_1$, one must construct the modulus $\delta : \N \to \N$ from the existential hypothesis \texttt{HasUnifDeriv A}:
\[
\operatorname{HasUnifDeriv}(A) \iff \exists f', \forall k, \exists N, \forall h, 0 < |h| < 2^{-N} \implies \left| \frac{f(x+h)-f(x)}{h} - f'(x) \right| < 2^{-k}
\]
In \texttt{QAnalysis.lean}, \texttt{A0.toA1} uses Lean's \texttt{Classical.choose} to extract $f'$ and $\delta$:
\begin{lstlisting}[language=Lean]
noncomputable def A0.toA1 (A : A0) (h : A0.HasUnifDeriv A) : A1 :=
  let f' := Classical.choose h
  let h_diff := Classical.choose_spec h
  let δ := fun k ↦ Classical.choose (h_diff k)
  { toA0 := A, f' := f', δ := δ, diff := ... }
\end{lstlisting}
\textbf{Measured Footprint:} \texttt{\#print axioms A0.toA1} produces \textbf{exactly \texttt{[Classical.choice]}}. Myhill~\cite{myhill1971} proved that there exists a computable $C^1$ function whose derivative is not computable; in our formalization, this computability barrier is reflected precisely as an inescapable dependency on choice.

\section{The Category $\Rep$ and Retract Equivalences}

We structure representations into a category $\Rep$:
\begin{definition}[Category $\Rep$]
An object $R \in \Rep(\alpha)$ is a carrier type with approximation relation $\operatorname{approx}$ and equivalence $\operatorname{equiv}$. A morphism $F : \operatorname{RepMorphism}(R_1, R_2)$ bundles a code transformer $f : C_1 \to C_2$ with a quantitative precision shift $\mu : \N \to \N$ such that:
\[
R_1.\operatorname{approx}(c, \mu(k), x) \implies R_2.\operatorname{approx}(f(c), k, F(x))
\]
\end{definition}

\begin{definition}[Representation Retract Preorder $\preceq$]
We define $R_1 \preceq R_2$ by the existence of a section-retraction pair:
\[
\iota : R_1 \to R_2, \quad \pi : R_2 \to R_1 \quad \text{such that} \quad \pi(\iota(c)) \sim c
\]
\end{definition}

We instantiate this preorder across the smoothness hierarchy by defining restricted subtypes with differentiability witnesses:
\begin{itemize}
    \item \texttt{RepA0diff} (Carrier: $\{ A : A_0 \mid \operatorname{HasUnifDeriv}(A) \}$)
    \item \texttt{RepE0diff} (Carrier: $\{ E : E_0 \mid \operatorname{HasSmoothDataE0}(E) \}$)
\end{itemize}

\begin{theorem}[Galois Equivalences]
In \texttt{HAomega/GaloisAdequacy.lean}:
\begin{lstlisting}[language=Lean]
theorem A1_le_A0diff (a b : Q) : (RepA1 a b) ⪯ (RepA0diff a b)
theorem A0diff_le_A1 (a b : Q) : (RepA0diff a b) ⪯ (RepA1 a b)
theorem A1_equiv_A0diff (a b : Q) : (RepA1 a b) ≃ᵣ (RepA0diff a b)

theorem E1_le_E0diff (a b : Q) : (RepE1 a b) ⪯ (RepE0diff a b)
theorem E0diff_le_E1 (a b : Q) : (RepE0diff a b) ⪯ (RepE1 a b)
theorem E1_equiv_E0diff (a b : Q) : (RepE1 a b) ≃ᵣ (RepE0diff a b)
\end{lstlisting}
\end{theorem}
$C^1$ data ($A_1$) and differentiable $C^0$ data ($A_0^{\text{diff}}$) represent the \textbf{same space up to Galois retract equivalence} ($A_1 \simeqR A_0^{\text{diff}}$).

\section{Limitative Results}

\begin{enumerate}
    \item \textbf{Friedman--Ko Complexity Barrier:} Friedman and Ko~\cite{friedman1982} proved that there exist polynomial-time computable $C^\infty$ functions whose definite integrals are $\#\mathbf{P}_1$-complete. Consequently, polynomial-time representation adequacy for $\text{EFTC2}$ cannot be achieved even with infinite smoothness data.
    \item \textbf{Construction Slack:} Our verified Riemann summation implementation \texttt{sqEx.intN} incurs a precision bound of $2^{2k+14}$ against an information-theoretic lower bound of $\approx 2^k$.
    \item \textbf{Baire Continuity vs. Metric Moduli:} The higher-type Baire-space continuity predicate \texttt{HasMod} measures oracle tape inspection, whereas analysis requires metric moduli on $\Q$. We record that metatheorems on higher-type functionals do not automatically yield metric moduli without explicit Kleene associates.
\end{enumerate}

\section{Related Work}

\begin{itemize}
    \item \textbf{C-CoRN (Cruz-Filipe et al.~\cite{cruzfilipe2004}):} Constructive analysis in Coq under uniform constructive assumptions. We deeply embed an object-level logic ($\HAomega$) and explicitly calibrate non-constructive gaps via kernel axiom auditing.
    \item \textbf{Incone (Steinberg, Th\'ery, Thies~\cite{steinberg2017,steinberg2022}):} Type-Two Effectivity over dialogue represented spaces. We connect represented spaces directly to formal proof extraction via \texttt{central\_adequacy\_theorem}.
    \item \textbf{Weihrauch Complexity (Brattka, Gherardi, Pauly~\cite{brattka2021}):} Orders multi-valued problems $f \le_W g$. Our retract relation $R_1 \preceq R_2$ orders the \textbf{representation spaces themselves}.
    \item \textbf{Proof Mining (Kohlenbach~\cite{kohlenbach2008}):} Extraction of moduli from non-constructive proofs using monotone Dialectica. We formalize modified realizability for pure constructive $\HAomega$.
\end{itemize}

\section{Conclusion}

By organizing representations into a structured hierarchy ($A_0 \subset A_1, E_0 \subset E_1$), formalizing the retract preorder $\preceq$, and using \texttt{\#print axioms} as a quantitative measurement instrument, we demonstrate that \textbf{the exact data requirements of classical theorems can be rigorously calibrated, mechanically checked, and connected to certified executable code.}

\bibliographystyle{plain}
\bibliography{references}

\end{document}
